\section{A smaller fixed support}\label{sec:modular-support}

We use a smaller support for the subsequent constructions. Replace the
support positions and exponent of Section~\ref{sec:input-unit}, keeping
its arithmetic structure. The resulting certificate still has 107
instructions. Fixed numerals are free in the complexity measure used
for the reductions below.

\subsection{Pair recovery over a finite field}

For $0\le i\le59$, let $r_i$ be the least nonnegative residue of $i^2$
modulo~61, and put
\[
 a_i=122i+r_i,\qquad v_i=1+3a_i.
\]
The sums $a_i+a_j$ distinguish unordered pairs, including repeated
indices. Indeed, $0\le r_i+r_j\le120<122$, so division by~122 recovers
the integer sum $s=i+j$ and the residue sum $r_i+r_j$. Over the field
with 61 elements these determine
\[
 i+j=s,\qquad ij=(s^2-r_i-r_j)/2.
\]
Thus the monic polynomial with roots $i,j$ is determined. Since the
indices $0,\ldots,59$ are distinct modulo~61, the integer pair is
determined as well. The division by~2 occurs in this proof of the
fixed support construction; it is not a certificate instruction.

Every $v_i$ is 1 modulo~3. Hence zero, a single positive weight, and
a pair sum are distinguished modulo~3. Within each class the weights
are distinct by the preceding argument. There are consequently
$1+60+1830=1891$ distinct homogeneous quadratic monomial weights in
the 61 coordinates. As before, the input has weight zero, $\delta$
has weight $v_0=1$, the original witnesses have weights $v_1,\ldots,v_{58}$,
and $u$ has weight $v_{59}$.

With $m=1832$ target rows, use the same formulas for their spacing:
\[
\begin{gathered}
 v_*=21607,\qquad D_*=4v_*+1=86429,\\
 T_*=(m+1)D_*+2v_*=158467571,\qquad t_j=T_*+jD_*,\\
 t_*=316719070,\qquad K=t_*+1=316719071,\qquad L=2^{32}.
\end{gathered}
\]
The exact margins are
\[
\begin{gathered}
 D_*>2v_*,\qquad T_*-2v_*>v_*,\qquad 2T_*-2v_*>t_*,\\
 3K+2=950157215<2^{30}<L.
\end{gathered}
\]
Different row intervals $[t_j-2v_*,t_j]$ are disjoint. A product of
a coefficient term and an ordinary quadratic monomial lies within
$2v_*$ of its row target, so no other row contributes there. Coefficient
terms lie above every low variable position, and products involving a
dummy row coordinate lie above every target. These are all the support
facts used in the preceding decoding proof.

Retain twice each ordinary homogeneous residual and the special
coefficient $\delta^2$. In particular, its coefficient 1 is placed at
$t_{m-1}-2$. Recompute $D,\ell_0,e_0,V$ with the new positions, choose
the power of two $z\ge2$ above the absolute coefficients of $D$, put
$Z=2z$, and choose a power of two
\[
 H>\max\{2Z^{2L+1},\ 4^{t_*+3}Z\,1900^2,\ 3L,\ 16\}.
\]
The number of coordinates remains 1893, so their coefficient sum remains
less than $1900b$. All bounds in the input-unit and positive-coefficient
proofs are therefore preserved. In particular, $e,\mathcal C<B^K$ and
$L>3K+2$ still give $2e\mathcal C^2<q$ after exponent recovery, and
evaluation at~4 still rules out the high-mask quotient ambiguity.
The fixed indices $Z,V,H$ are computed from the represented set,
independently of the queried input.

\subsection{The even exponent and the preserved certificate}

No preceding Pell argument requires $L$ to be odd. The signed congruence
uses the odd index $J=2r+1$. The other index argument uses only
$0<L<J<a$ and the congruence $\psi_a(t)\equiv t\pmod{a-1}$, valid for
every $t$. These hypotheses still hold because $3L<B\le n\le r$.
Necessity gives the positive integer
$\Delta=(\psi_a(L)-L)/(a-1)$ since $L>1$, and the positive gap follows
from $J>L$. The exponent comparison
$B^{3L}\le B^B\le n^n\le U^r<a$ is unchanged. Thus $q=B^L$ is
recovered, and every decoding and necessity argument still applies.

\begin{theorem}\label{thm:modular107}
The modular support gives a membership-equivalent system in 34 positive
unknowns and 22 equations with \textbf{107 arithmetic instructions}:
59 multiplications and 48 additions.
\end{theorem}

\begin{proof}
The support and exponent arguments establish the same encoding theorem.
The explicit schedule moves the existing register $a_-=a-1$ before E14 and
replace its computations of $4a$ and $4a-5$ by
\[
 a_4=4a_-,\qquad a_{4-5}=a_4-1.
\]
The latter is identically $4a-5$, so all 107 arithmetic instructions
are preserved. The complete schedule is serialized and checked by
\file{round13\_1980\_certificate.py}. It checks every primitive and
source residual identity and all 1891 monomial weights and support
margins. Its 59 multiplications and 48 additions have exactly the
same costs as in Appendix~\ref{app:107}.
\end{proof}
